\documentclass[12pt]{article}

% ------------------------------------------------------------
% Pacotes básicos
% ------------------------------------------------------------
\usepackage[T1]{fontenc}
\usepackage[utf8]{inputenc}
%\usepackage[brazil]{babel}
\usepackage[english]{babel}

\usepackage{lipsum}         % texto exemplo
\usepackage{graphicx}       % imagens
\usepackage{float}          % controle de posição de figuras
\usepackage{booktabs}       % tabelas
\usepackage{amsmath,amssymb}% matemática
\usepackage{todonotes}
\usepackage{hyperref}       % links clicáveis
\usepackage{caption}        % legendas personalizadas
\hypersetup{
  colorlinks   = true,
  urlcolor     = blue,
  linkcolor    = blue,
  citecolor    = blue
}
\usepackage{tikz}            % diagramas
\usepackage{microtype}       % melhorias tipográficas
\usepackage{enumitem}
\usetikzlibrary{patterns.meta, arrows.meta, positioning, shapes.geometric}

% ------------------------------------------------------------
% Todos and notes
% ------------------------------------------------------------
%\setlength{\marginparwidth}{2cm}
%\usepackage[colorinlistoftodos,textsize=scriptsize]{todonotes}

% ------------------------------------------------------------
% Bibliografia e Citações
% ------------------------------------------------------------
\usepackage{csquotes} % citações
\usepackage[backend=biber, style=abnt]{biblatex}
\addbibresource{references.bib}

% ------------------------------------------------------------
% Configurações do documento
% ------------------------------------------------------------
% Page setup
\usepackage{geometry}
\geometry{a4paper, margin=2.5cm}

% Paragraph spacing
\usepackage{parskip}

% Configuração dos títulos dos apêndices
\usepackage{titlesec}
\renewcommand{\appendixname}{APÊNDICE}

% ------------------------------------------------------------
% Variáveis do Documento
% ------------------------------------------------------------

\newcommand{\instituicao}{UNIVERSIDADE FEDERAL DE SANTA CATARINA}
\newcommand{\departamento}{Programa de Pós-Graduação em Economia}
\newcommand{\curso}{Mestrado}
\newcommand{\disciplina}{Métodos Quantitativos}
\newcommand{\titulo}{Sundaram}
\newcommand{\subtitulo}{Exercícios - Capítulo 1}
\newcommand{\autor}{}
%\newcommand{\orientador}{}
\newcommand{\local}{Florianópolis}
\newcommand{\dataentrega}{\today}
%\newcommand{\resumo}{
%\lipsum[1][1-8] % Resumo do trabalho
%}

% ------------------------------------------------------------
% Título simples
% ------------------------------------------------------------
\newcommand{\imprimirtitulo}{
\noindent

\begin{center}
  {\Large\bfseries\MakeUppercase{\instituicao}\par}
  \vspace{0.2em}
  {\large \departamento\par}
  \vspace{0.2em}
  {\normalsize\textbf{\disciplina}\ (\curso)\par}
  \vspace{1.4em}
  {\LARGE\bfseries \titulo\par}
  \vspace{0.25em}
  {\large\itshape \subtitulo\par}
  \vspace{1.5em}
\end{center}

\noindent
\begin{minipage}{0.48\textwidth}
  \raggedright\normalsize \autor
\end{minipage}\hfill
\begin{minipage}{0.48\textwidth}
  \raggedleft\normalsize \local, \dataentrega
\end{minipage}

\vspace{1.0em}
\hrule
}


% ------------------------------------------------------------
% Capa
% ------------------------------------------------------------

\begin{document}
\imprimirtitulo

% ----------------------------------------------------------------------------
% Sumário
% ----------------------------------------------------------------------------

\renewcommand{\contentsname}{Sumário}
\tableofcontents
\vspace{1.5em}
\hrule

% ----------------------------------------------------------------------------
% Corpo do documento
% ----------------------------------------------------------------------------

Show that equality obtains in the Cauchy--Schwartz inequality if and only if the vectors $x$ and $y$ are collinear (i.e., either $x = ay$ for some $a \in \mathbb{R}$ or $y = bx$ for some $b \in \mathbb{R}$). (Note: This is a difficult question. Prove the result for unit vectors first, i.e., for vectors $x$ and $y$ which satisfy $x \cdot x = y \cdot y = 1$. Then reduce the general case to this one.)

\newpage

% Autor - Gabriel Cintra
\section{Resolução 1}
  
  \textbf{Answer:} 

  \textbf{Claim:} 
  Let $x$ and $y$ be two vectors in $\mathbb{R}^n$. 

  Then $|x \cdot y| = (x \cdot x)^{1/2}(y \cdot y)^{1/2} \Leftrightarrow (\exists a \in \mathbb{R})(x = ay) \lor (\exists b \in \mathbb{R})(y = bx).$

  \textbf{Proof:}

  The proof will be structured in three steps.

  We begin by proving $(\Rightarrow)$, using two different approaches. 
  
  We then prove $(\Leftarrow)$ to complete the proof.

  We assume $x \ne 0$ since in the trivial case where $x=0$ the equality holds and $x,y$ are collinear.

  \textbf{Step 1: $(\Rightarrow)$: Approach 1: By algebraic manipulation:}

  From the definition of Euclidean inner product we can express $|x \cdot y| = (x \cdot x)^{1/2}(y \cdot y)^{1/2}$ in summation notation, as follows:
  
  $\left| \sum_{i=1}^n x_i y_i \right|=\left( \sum_{i=1}^n x_i^2 \right)^{1/2}\left( \sum_{i=1}^n y_i^2 \right)^{1/2}$,

  $\left(\left( \sum_{i=1}^n x_i y_i \right)^2\right)^{1/2}=\left( \sum_{i=1}^n x_i^2 \right)^{1/2}\left( \sum_{i=1}^n y_i^2 \right)^{1/2}$,

  $\left( \sum_{i=1}^n x_i y_i \right)^2= \sum_{i=1}^n x_i^2 \cdot \sum_{i=1}^n y_i^2 $,

  $(x_1y_1 + x_2y_2 + \cdots + x_ny_n)^2=(x_1^2 + x_2^2 + \cdots + x_n^2)(y_1^2 + y_2^2 + \cdots + y_n^2)$

  $\sum_{i=1}^n x_i^2 y_i^2 + 2\sum_{1 \le i < j \le n} x_i x_j y_i y_j = \sum_{i=1}^n \sum_{j=1}^n x_i^2 y_j^2$

  $\sum_{i=1}^n x_i^2 y_i^2 + 2\sum_{1 \le i < j \le n} x_i x_j y_i y_j = \sum_{i=1}^n x_i^2 y_i^2 + \sum_{1 \le i < j \le n} (x_i^2 y_j^2 + x_j^2 y_i^2)$

  $2\sum_{1 \le i < j \le n} x_i x_j y_i y_j = \sum_{1 \le i < j \le n} (x_i^2 y_j^2 + x_j^2 y_i^2)$

  $\sum_{1 \le i < j \le n} (x_i^2 y_j^2 + x_j^2 y_i^2 - 2x_i x_j y_i y_j)=0$

  $\sum_{1 \le i < j \le n} (x_i y_j - x_j y_i)^2 = 0$.

  A sum of squares can only be zero if each term is equal to zero. Hence

  $x_i y_j = x_j y_i \qquad \text{for all } 1 \le i < j \le n.$

  If $y=0$, then $y = 0x$, hence $x$ and $y$ are collinear.

  If $y \ne 0$, choose $k$ such that $y_k \ne 0$, and let $a = \dfrac{x_k}{y_k}$.

  Then for every $i$ we have $x_i y_k = x_k y_i$, so $x_i = \dfrac{x_k}{y_k} y_i = ay_i$.

  Therefore $x = ay$, and hence $x$ and $y$ are collinear.

  \newpage
  \textbf{Step 2: $(\Rightarrow)$: Approach 2: By vector Decomposition:}

  From the definition of Euclidean inner product we can express $|x \cdot y| = (x \cdot x)^{1/2}(y \cdot y)^{1/2}$ as follows:

  $(x'y)^2=x'xy'y$

  $y'y = \frac{(x'y)^2}{x'x}$

  We then decompose $y$ onto a component parallel to $x$ and another component orthogonal to it, thus, for some $c \in \mathbb{R}$:

  $y = cx + e$, where $x'e = 0$, hence

  $x'y = cx'x + x'e$, since $x'e = 0$,

  $x'y = cx'x$, and,

  $c = \frac{x'y}{x'x}$.

  Therefore we can express $y'y$ as:

  $y'y=(cx+e)'(cx+e)$

  $y'y=(cx'+e')(cx+e)$

  $y'y=cx'cx+e'cx+cx'e+e'e$

  $y'y=cx'cx+e'e$

  $y'y=c^2x'x+e'e$

  $y'y=\frac{(x'y)^2}{x'x}+e'e$.

  Hence $e'e=0$, which can only be true if $e=0$.

  Therefore $y = cx$.

  \textbf{Step 3: $(\Leftarrow)$:}

  Assume that $x,y$ are collinear for some $c$, hence:

  $y = cx$

  $y'y=cy'x$

  $y'y=cx'y$.

  Choose $c = \frac{x'y}{x'x}$, then:

  $y'y=\frac{x'y}{x'x}x'y$

  $(x'y)^2=x'xy'y$, taking the square root in both sides and using the inner product notation:

  $|x \cdot y| = (x \cdot x)^{1/2}(y \cdot y)^{1/2}$.

  $\square$

% ----------------------------------------------------------------------------
% Referências
% ----------------------------------------------------------------------------

%\newpage
%\renewcommand{\refname}{REFERÊNCIAS}
%\addcontentsline{toc}{section}{REFERÊNCIAS}
%\printbibliography

% ----------------------------------------------------------------------------
% Apêndice
% ----------------------------------------------------------------------------

%\newpage
%\appendix
%\section*{APÊNDICES} % título de abertura
%\addcontentsline{toc}{section}{APÊNDICES} % entra no sumário

% ------------------------------
% APÊNCICE A
% ------------------------------

%\section{Referência dos dados utilizados}
%\subsection{Produto}

% ------------------------------
% Final do documento
% ------------------------------

\end{document}
